% Requires amsmath, amssymb, amsthm, booktabs, xurl.
% This is a source-level alternative, not a claim of a 194-state realization.
\section{An arbitrary-consequent gate and a ten-entry four-field interface}
\label{sec:gate-x}

This section gives a different numerical calculus and an exact source program
for it.  The source has 3,666 primitive sites and 4,772 branch outcomes.
Its first binary realization has 198 working states; a smaller grammar is a
separate optimization problem.  Thus the construction below is not counted as
a solution of the requested 194-state bound.  It records a mathematically
validated semantic alternative, including the source arithmetic and the
all-store proof compiler, so that future physical compression can reuse it.

\subsection{A conservative change to the numerical calculus}

Retain Cantor pairing and the original numerical constructors
\[
 p(a,b)=\frac{(a+b)(a+b+1)}2+b,\qquad I(a,b)=4p(a,b),
\]
\[
 M(a,b)=I(a,b)+1,\quad U(a,B)=I(B,a)+2,\quad E(a,b)=I(a,b)+3.
\]
Write $A\to B=I(A,B)$, $N(A)=A\to1$,
$A\wedge B=N(A\to N(B))$, and $\exists_a B=N(U(a,N(B)))$.
The distinguished contradiction is 1; the truth code 0 satisfies
$I(0,0)=0$ literally.  The six ordinary axiom heads are
\begin{align*}
 h_1(D,H)&=(\exists_D U(D,H))\to H,\\
 h_2(D,A,H)&=U(D,U(A,H))\to U(A,U(D,H)),\\
 h_3(D,A,H)&=M(H,D)\to(E(H,A)\to M(A,D)),\\
 h_4(D,A,H)&=E(H,D)\to(M(A,H)\to M(A,D)),\\
 h_5(D,H)&=\exists_D E(H,D),\\
 h_6(D,A,H)&=U(D,A)\to(U(D,A\to H)\to U(D,H)).
\end{align*}
The additional unchanged heads are
\[
 \begin{aligned}
 \Gamma(B)&=B\to U(B,B),\\
 \rho_8(D,A,H,u)&=((H\to u)\to D)\to((D\to H)\to(A\to H)),
 \qquad u\le D.
\end{aligned}
\]
The inherited restricted implicational basis uniformly derives $K$, $S$,
identity, composition, and deduction by primitive recursive proof recipes.
No gate axiom is used in those recipes.

For $c\ne a,b$, put
\begin{align*}
 V(a,b,c)&=M(a,b)\to U(c,M(a,b)),\\
 P(a,b,c)&=c\to M(c,b),\\
 S(a,b,c)&=U\bigl(b,(M(a,b)\to P(a,b,c))\wedge(M(c,b)\to c)\bigr),\\
 Q(a,b,c)&=U(c,S(a,b,c)\to N(P(a,b,c))),\\
 B(a,b,c)&=V(a,b,c)\to Q(a,b,c).
\end{align*}
Replace the old gate $G(a,b,c)=B(a,b,c)\to1$ by the family
\[
                 G_z(a,b,c)=B(a,b,c)\to z,
                 \qquad z\in\mathbb N.                  \tag{GX1}
\]
Keep both inference rules, modus ponens and theorem generalization.
Call the resulting numerical calculus $\mathcal C_X$; call the original
one $\mathcal C_4$.  The failed-guard output remains 0, whose fixed proof is
$\rho_8(0,0,0,0)$.

\paragraph{Interderivability lemma.}
There are uniform primitive recursive proof translations in both directions
between $\mathcal C_4$ and $\mathcal C_X$, preserving the exact numerical
conclusion.  Their correctness is provable in iPRA.
\begin{proof}
The original gate is the $z=1$ instance of (GX1), so one translation merely
changes the gate-leaf tag and adds the parameter 1.  For the converse, first
derive explosion using only $h_1$ and propositional logic.  Put
$F=U(0,N(U(0,z)))$.  Then
\[
 h_1(0,z)=(F\to1)\to z,\qquad K(1,F)=1\to(F\to1).
\]
Their composition proves $1\to z$.  A second composition with
$G(a,b,c)=B(a,b,c)\to1$ proves $G_z(a,b,c)$.
Every displayed derivation is a fixed finite $K,S$, MP recipe parameterized
by $a,b,c,z$.  In particular it neither invokes semantic validity nor depends
on a decision procedure for numerical theoremhood.  The unchanged guard
$c\ne a,b$ licenses the old gate used in that recipe.

Translate a finite proof list in order, replacing only its gate rows and
relocating earlier indices after concatenation.  The finite recipes and
index operations are primitive recursive.  A bounded Boolean invariant
checks every translated row and its conclusion; quantifier-free induction
on the input-list length proves the two translation theorems.  This also
handles the numerical zero alias, since the constructors in the recipes
are interpreted by their literal numerical equations.
\end{proof}

Consequently iPRA proves equality of the two numerical refutability
predicates.  The change is conservative even though its emitter accepts a
new independent formula parameter.

\subsection{An exact ten-entry AX command}

Use the store $(t,w,d,h)$ and the total Cantor stack notation
$[a;w]=p(a,w)$.  Its projections satisfy
$\operatorname{hd}[a;w]=a$ and $\operatorname{tl}[a;w]=w$, including
$[0;0]=0$.  Set $D=d+x$, $H=h+y$, $A=c$ and define
\[
\begin{array}{c|l}
 j&q_j\\\hline
 0&G_z(H,D,A),\ \text{or }0\text{ if }A=D\text{ or }A=H\\
 1&\Gamma(A)\\
 2&h_6(D,A,H)\\
 3&h_5(D,H)\\
 4&h_4(D,A,H)\\
 5&h_3(D,A,H)\\
 6&h_2(D,A,H)\\
 7&h_1(D,H)\\
 8&\rho_8(D,A,H,d)\\
 9&U(d,h).
\end{array}
\]
The proposed AX interface is
\[
 \mathsf A_{c,x,y,z}(t,w,d,h)
       =(q_0,[q_1,\ldots,q_9,t;w],0,0).                 \tag{GX2}
\]
It has four actual fields.  The fourth field becomes the gate consequent;
it does not become the offset and does not cause a fallthrough MP.
All $q_0,\ldots,q_8$ are congruent to 0 modulo 4, and $q_9$ is congruent
to 2 modulo 4 and at least 2.  Thus none is 1, for any fourth field.

Retain the cached-only Shift and the no-push MP:
\[
 \mathsf S(t,w,d,h)=
 \begin{cases}\mathsf{halt},&h=1,\\
 (h,\operatorname{tl}w,d,\operatorname{hd}w),&h\ne1,
 \end{cases}                                           \tag{GX3}
\]
\[
 \mathsf{MP}_x(t,w,d,h)=
 \begin{cases}(b,w,b,0),&I(t,b)=h,\\
 (0,w,b,0),&I(t,b)\ne h,
 \end{cases}\qquad b=d+x.                              \tag{GX4}
\]
Only Shift can halt.

\subsection{Source realization and exact control flow}

The source register order is $(N,T,W,A,D,H,X,Y,L)$.
At AX entry, the last register is $p(y,p(x,p(z,\ell)))$.
The unchanged prefix constructs the nine stack heads, the old gate's
antecedent $B(H,D,A)$, and retains the guard parameters $A,D,H$.
In the final negation macro, replace its single increment of $T$ by
the two existing arithmetic macros
\[
             \mathcal U_L\ ;\ (X\longrightarrow T).     \tag{GX5}
\]
Here $\mathcal U_L$ unpairs $L$ into $X$ and $L$ and restores $Y=0$.
The preceding Push gives $T=X=Y=0$; thus (GX5) changes
\[
 (T,X,Y,L)=(0,0,0,p(z,\ell))
            \longmapsto(z,0,0,\ell).
\]
It leaves $A,D,H$ intact.  The following Cons pops $B(H,D,A)$ and
constructs $I(B(H,D,A),z)$.  The existing CopyGuard and CompareGuard
then apply both original tests and clear the guard and scratch registers.
An ordinary BREAK ends AX without pushing its visible result.

The tree token replacement is
\[
 \text{leaf ordinal }185:\qquad
  (\texttt{leaf},\texttt{inc},1)
      \longmapsto(\texttt{node32},\texttt{node42}).
\]
Parenthesizing that pair supplies the existing native macro entry contexts:
$\texttt{node32}$ enters native control 104 after a 0; $\texttt{node42}$
enters native control 81 after a 1.  These entry contexts are physical
requirements, rather than consequences of the arithmetic macro equations.

The independent checker \path{check_gate_x.py} assembles every primitive
template and compares the complete graph with a second resolution of the
source tree.  It establishes all 3,666 instruction rows and their outgoing
edges, with the following ports:
\[
\begin{array}{c|rrrrrrrr}
 &\mathrm{fetch}&\mathrm{decode}&\mathrm{AX}&\mathrm{AXbreak}
 &\mathrm{MP}&\mathrm{MPcompare}&\mathrm{return}&\mathrm{Shift}\\\hline
 \mathrm{site}&25&34&50&3566&3567&3637&3654&3655.
\end{array}
\]
In particular site 3566 jumps to 3654, which jumps to 25.  The path skips
the MP block beginning at 3567.  Independent free-algebra substitution
through the complete AX word verifies the exact nine-register endpoint
\[
 (n,q_0,[q_1,\ldots,q_9,t;w],0,0,0,0,0,\ell).         \tag{GX6}
\]
Both equality guards remain symbolic, and every intermediate scratch
precondition is checked.  This is an expression identity using pairing
projections, rather than a test on finitely many formula codes.

The inherited written arithmetic-template proofs supply exact primitive
recursive clocks.  The new two-macro clock is
\[
           K_{\mathcal U}(p(z,\ell))+3z+1,
\]
where $K_{\mathcal U}$ is the inherited total unpair clock.  The terminal
AX break and return together add two primitive steps.  Composing the
fixed macro clocks at their successive arithmetic endpoints proves (GX6)
at an explicit PR duration.  The Shift, MP, fetch and decode arguments are
unchanged except for their finite site relocation.  Consequently the
source-to-command simulation is a uniform iPRA proof once those inherited
template lemmas are used.  A binary realization additionally requires its
own literal contracts, frame closure and complete blank initialization.

\subsection{All-store instruction words}

Commands execute from left to right.  Put $A_0=\mathsf A_{0,0,0,0}$.
From (GX2)--(GX3), for every store and all four fields,
\begin{align*}
 \mathsf A_{c,x,y,z};\mathsf S^9(t,w,d,h)
       &=(q_8,[t;w],0,q_9),\\
 \mathsf A_{c,x,y,z};\mathsf S^{10}(t,w,d,h)
       &=(q_9,w,0,t).
\end{align*}
These executions never test the saved top; their tested generated values
are different from 1.  Define
\[
 \mathsf{Push}=A_0;\mathsf S^9;\mathsf{MP}_0,\qquad
 \mathsf{Pop}=A_0;\mathsf S^{11},\qquad
 \mathsf{Discard}=\mathsf{MP}_0,\qquad
 \mathsf{Offset}_i=\mathsf{MP}_i.
\]
Their exact contracts are
\begin{align*}
 \mathsf{Push}(t,w,d,h)&=(0,[t;w],0,0),\\
 \mathsf{Pop}(t,[a;w],d,h)
   &=\begin{cases}\mathsf{halt},&t=1,\\(t,w,0,a),&t\ne1,\end{cases}\\
 \mathsf{Discard}(t,w,0,h)&=(0,w,0,0),\\
 \mathsf{Offset}_i(0,w,0,0)&=(0,w,i,0).
\end{align*}
The Push comparison fails by residues: an implication cannot equal $q_9$.
The eleventh Shift in Pop tests the saved top.  Discard's two branches
have the same zero result, and $I(0,i)>0$ for $i>0$ proves Offset.
All statements include the total stack's zero aliases.

For the gate $q_0$, the selector is
\[
        \mathsf{Discard};\mathsf A_{A,D,H,z};\mathsf{Pop}^{10}.
                                                               \tag{GX7}
\]
For $1\le j\le8$, the selectors are
\[
 \mathsf{Discard};\mathsf A_{A,D,H,z};
                    \mathsf S^{j+1};\mathsf{Pop}^{9-j}.        \tag{GX8}
\]
From every $(t,w,0,0)$, these words return $(q_j,w,0,0)$.
For (GX7) the ten popped entries are $q_1,\ldots,q_9,0$; for (GX8)
the remaining entries are $q_{j+2},\ldots,q_9,0$.
Pop preserves the selected $q_j\ne1$, and its last removed entry is 0,
so its final cache is 0.  The restricted head with $u\le D$ is selected by
\[
 \mathsf{Discard};\mathsf{Offset}_u;
 \mathsf A_{A,D\mathbin{\dot-}u,H,0};\mathsf S^9;\mathsf{Pop}.
                                                               \tag{GX9}
\]
The old offset is exactly $u$, the accumulated parameter is $D$, and AX
itself clears the offset.  The generalization word is
\[
 \mathsf{Gen}_i=\mathsf{Push};\mathsf{Offset}_i;\mathsf S;
                  \mathsf A_{i,0,0,0};\mathsf S^{10}.           \tag{GX10}
\]
It maps $(a,w,0,0)$ to $(U(i,a),w,0,0)$ even when $a=1$:
the only Shift before AX sees cache 0, and AX clears the cache containing
$a$ before another Shift tests it.

From $(0,[a,I(a,b),0;w],0,0)$ the suffix
\[
                    \mathsf S;\mathsf S;\mathsf{MP}_b;\mathsf{Pop}
\]
halts if $a=1$ or $b=1$ and otherwise returns $(b,w,0,0)$.
As in AX11H, the reserved zero protects the base stack.  Compile an MP
node, with major proof $e_1$ and minor proof $e_0$, by
\[
 W_e=\mathsf{Discard};\mathsf{Push};W_{e_1};\mathsf{Push};
 W_{e_0};\mathsf{Push};\mathsf S;\mathsf S;\mathsf{MP}_b;\mathsf{Pop}.
\]
Axiom rows use (GX7)--(GX9); generalization appends (GX10).  The exact
compiler invariant is
\[
 W_e(t,w,0,0)=\mathsf{halt}\quad\text{or}\quad
 W_e(t,w,0,0)=(\operatorname{concl}(e),w,0,0).            \tag{GX11}
\]
Therefore $A_0;W_e;\mathsf{Pop}$ halts from every store when $e$ is a
refutation.  Its prefix normalizes the offset and cache without assuming
that the incoming store is empty.

\paragraph{Induction strength.}
For completeness, the usual phrase ``induct on proofs for every base
store'' can be replaced by a quantifier-free induction.  Encode a proof
as $N$ ordered lines; let $L$ bound all its compiled word lengths and $K$
all command codes.  Define the monotone PR function
\[
 \Phi_K(s)=1+\max\bigl(s,\max_{c\le K,\ u\le s}\operatorname{next}(c,u)\bigr).
\]
For an arbitrary input bound $s_0$, set
$B(j)=\Phi_K^{L(N-1-j)}(s_0)$.  Prove successively that (GX11) holds
for every earlier line $i$ and every normalized store code at most
$B(i)$.  This is a bounded Boolean conjunction, hence an allowed
quantifier-free induction predicate.  Every recursive premise call from
line $j$ is preceded by at most $L$ commands and has input at most
$\Phi_K^L(B(j))\le B(i)$ for $i<j$.  The local contracts above complete
the induction step.  At $j=N-1$ the original arbitrary bound is $s_0$;
universal generalization yields (GX11).  No induction over an unbounded
universal environment predicate is needed.

\subsection{Scheduler completeness and soundness}

Use the same actual field order as AX11H:
\[
 e(\epsilon)=0,\quad e(c::w)=1+b_c(e(w)),
\]
\[
 b_{\mathsf S}(r)=p(0,r),\quad
 b_{\mathsf{MP}_x}(r)=p(1,p(x,r)),\quad
 b_{\mathsf A_{c,x,y,z}}(r)=p(c+2,p(y,p(x,p(z,r)))).       \tag{GX12}
\]
The numerical meaning of the fourth field changed; its position did not.
Fetch uses raw seed $n$ and increments $n$ when the continuation is zero;
otherwise it uses continuation minus one and leaves $n$ unchanged.
Every residual continuation is a nested Cantor component of that raw
seed, so each positive continuation strictly decreases.  Candidate $n$
therefore drains in at most $n+1$ commands unless it halts.  A candidate
$N\ge2$ is reached within the loose PR bound $(N+1)^2$ from the initial
next-unused counter 2, unless a prior halt has already occurred.

Prefix the all-store refutation word by $\mathsf A_{k,0,0,0}$.
Its raw seed is
\[
     p(k+2,p(0,p(0,p(0,e(\text{refutation suffix}))))).
\]
These seeds are cofinal, since they are at least $k+2$; the suffix works
for every carried store.  Thus the scheduler halts from its clean start
whenever $\mathcal C_X$ has a finite refutation.  The projection equations,
draining bound and fixed clocks give explicit primitive recursive witnesses
and quantifier-free induction proofs of this direction.

Conversely, annotate the visible top, cache and canonical Cantor-stack
entries by indices in a finite valid $\mathcal C_X$ proof list.  The virtual
zero entries of numerical stack 0 use the fixed zero proof.  AX appends
the head proofs in (GX2); $q_9$ is generalized from the old cache,
$q_8$ obeys its actual guard $d\le d+x$, and $q_0$ is its new gate axiom
or zero.  The saved top already has a proof.  Successful MP combines the
two indicated proofs; failed MP produces zero.  Shift moves annotations,
and a halting Shift exposes an annotated proof of 1 in its cache.
The annotation updates and their finite validity checker are PR, so
quantifier-free induction on command time extracts a numerical refutation
from every scheduler halt.

Writing $R_X$ for this exact scheduler, we have an ordinary written iPRA
proof of
\[
 \operatorname{Halt}(R_X)
  \longleftrightarrow \exists e\,\operatorname{Prf}_{\mathcal C_X}(e,1)
  \longleftrightarrow \exists e\,\operatorname{Prf}_{\mathcal C_4}(e,1).
                                                               \tag{GX13}
\]
Together with the source binding and the inherited PA calibration, this
establishes the intended source-level nonhalting/consistency equivalence.
A claim about a specified literal binary table still requires successful
physical certificates for that exact table.  Neither (GX13) nor the source
site reduction alone establishes a smaller binary state count.
